\section{Gysin homomorphisms and rational equivalence}
\label{section-gysin}

\noindent
In this section we use the key formula to show the Gysin homomorphism
factor through rational equivalence. We also prove an important
commutativity property.

\begin{lemma}
\label{lemma-gysin-factors-general}
Let $(S, \delta)$ be as in Situation \ref{situation-setup}.
Let $X$ be locally of finite type over $S$.
Let $X$ be integral and $n = \dim_\delta(X)$.
Let $i : D \to X$ be an effective Cartier divisor.
Let $\mathcal{N}$ be an invertible $\mathcal{O}_X$-module
and let $t$ be a nonzero meromorphic section of $\mathcal{N}$.
Then $i^*\text{div}_\mathcal{N}(t) = c_1(\mathcal{N}|_D) \cap [D]_{n - 1}$
in $\CH_{n - 2}(D)$.
\end{lemma}

\begin{proof}
Write $\text{div}_\mathcal{N}(t) = \sum \text{ord}_{Z_i, \mathcal{N}}(t)[Z_i]$
for some integral closed subschemes $Z_i \subset X$ of $\delta$-dimension
$n - 1$. We may assume that the family $\{Z_i\}$ is locally
finite, that $t \in \Gamma(U, \mathcal{N}|_U)$ is a generator
where $U = X \setminus \bigcup Z_i$, and that every irreducible component
of $D$ is one of the $Z_i$, see
Divisors, Lemmas \ref{divisors-lemma-components-locally-finite},
\ref{divisors-lemma-divisor-locally-finite}, and
\ref{divisors-lemma-divisor-meromorphic-locally-finite}.

\medskip\noindent
Set $\mathcal{L} = \mathcal{O}_X(D)$. Denote
$s \in \Gamma(X, \mathcal{O}_X(D)) = \Gamma(X, \mathcal{L})$
the canonical section. We will apply the discussion of
Section \ref{section-key} to our current situation.
For each $i$ let $\xi_i \in Z_i$ be its generic point. Let
$B_i = \mathcal{O}_{X, \xi_i}$. For each $i$ we pick generators
$s_i \in \mathcal{L}_{\xi_i}$ and $t_i \in \mathcal{N}_{\xi_i}$
over $B_i$ but we insist that we pick $s_i = s$ if $Z_i \not \subset D$.
Write $s = f_i s_i$ and $t = g_i t_i$ with $f_i, g_i \in B_i$.
Then $\text{ord}_{Z_i, \mathcal{N}}(t) = \text{ord}_{B_i}(g_i)$.
On the other hand, we have $f_i \in B_i$ and
$$
[D]_{n - 1} = \sum \text{ord}_{B_i}(f_i)[Z_i]
$$
because of our choices of $s_i$. We claim that
$$
i^*\text{div}_\mathcal{N}(t) =
\sum \text{ord}_{B_i}(g_i) \text{div}_{\mathcal{L}|_{Z_i}}(s_i|_{Z_i})
$$
as cycles. More precisely, the right hand side is a cycle
representing the left hand side. Namely, this is clear by our
formula for $\text{div}_\mathcal{N}(t)$ and the fact that
$\text{div}_{\mathcal{L}|_{Z_i}}(s_i|_{Z_i}) = [Z(s_i|_{Z_i})]_{n - 2} =
[Z_i \cap D]_{n - 2}$ when $Z_i \not \subset D$ because in
that case $s_i|_{Z_i} = s|_{Z_i}$ is a regular section, see
Lemma \ref{lemma-compute-c1}. Similarly,
$$
c_1(\mathcal{N}) \cap [D]_{n - 1} =
\sum \text{ord}_{B_i}(f_i) \text{div}_{\mathcal{N}|_{Z_i}}(t_i|_{Z_i})
$$
The key formula (Lemma \ref{lemma-key-formula}) gives the equality
$$
\sum \left(
\text{ord}_{B_i}(f_i) \text{div}_{\mathcal{N}|_{Z_i}}(t_i|_{Z_i}) -
\text{ord}_{B_i}(g_i) \text{div}_{\mathcal{L}|_{Z_i}}(s_i|_{Z_i}) \right) =
\sum \text{div}_{Z_i}(\partial_{B_i}(f_i, g_i))
$$
of cycles. If $Z_i \not \subset D$, then $f_i = 1$ and hence
$\text{div}_{Z_i}(\partial_{B_i}(f_i, g_i)) = 0$. Thus we get a rational
equivalence between our specific cycles representing
$i^*\text{div}_\mathcal{N}(t)$ and $c_1(\mathcal{N}) \cap [D]_{n - 1}$
on $D$. This finishes the proof.
\end{proof}

\begin{lemma}
\label{lemma-gysin-factors}
Let $(S, \delta)$ be as in Situation \ref{situation-setup}.
Let $X$ be locally of finite type over $S$.
Let $(\mathcal{L}, s, i : D \to X)$ be as in
Definition \ref{definition-gysin-homomorphism}.
The Gysin homomorphism factors through rational equivalence to
give a map $i^* : \CH_{k + 1}(X) \to \CH_k(D)$.
\end{lemma}

\begin{proof}
Let $\alpha \in Z_{k + 1}(X)$ and assume that $\alpha \sim_{rat} 0$.
This means there exists a locally finite collection of integral
closed subschemes $W_j \subset X$ of $\delta$-dimension $k + 2$
and $f_j \in R(W_j)^*$ such that
$\alpha = \sum i_{j, *}\text{div}_{W_j}(f_j)$.
Set $X' = \coprod W_i$ and consider the diagram
$$
\xymatrix{
D' \ar[d]_q \ar[r]_{i'} & X' \ar[d]^p \\
D \ar[r]^i & X
}
$$
of Remark \ref{remark-pullback-pairs}. Since $X' \to X$ is proper
we see that $i^*p_* = q_*(i')^*$ by Lemma \ref{lemma-closed-in-X-gysin}.
As we know that $q_*$ factors through rational equivalence
(Lemma \ref{lemma-proper-pushforward-rational-equivalence}), it suffices
to prove the result for $\alpha' = \sum \text{div}_{W_j}(f_j)$
on $X'$. Clearly this reduces us to the case where $X$ is integral
and $\alpha = \text{div}(f)$ for some $f \in R(X)^*$.

\medskip\noindent
Assume $X$ is integral and $\alpha = \text{div}(f)$ for some $f \in R(X)^*$.
If $X = D$, then we see that $i^*\alpha$ is equal
to $c_1(\mathcal{L}) \cap \alpha$.
This is rationally equivalent to zero by Lemma \ref{lemma-factors}.
If $D \not = X$, then we see that $i^*\text{div}_X(f)$ is equal to
$c_1(\mathcal{O}_D) \cap [D]_{n - 1}$ in $\CH_{n - 2}(D)$ by
Lemma \ref{lemma-gysin-factors-general}. Of course
capping with $c_1(\mathcal{O}_D)$ is the zero map
(Lemma \ref{lemma-c1-cap-additive}).
\end{proof}

\begin{lemma}
\label{lemma-gysin-back}
Let $(S, \delta)$ be as in Situation \ref{situation-setup}. Let $X$ be locally
of finite type over $S$. Let $(\mathcal{L}, s, i : D \to X)$ be as in
Definition \ref{definition-gysin-homomorphism}. Then
$i^*i_* : \CH_k(D) \to \CH_{k - 1}(D)$ sends $\alpha$ to
$c_1(\mathcal{L}|_D) \cap \alpha$.
\end{lemma}

\begin{proof}
This is immediate from the definition of $i_*$ on cycles
and the definition of $i^*$ given in
Definition \ref{definition-gysin-homomorphism}.
\end{proof}

\begin{lemma}
\label{lemma-gysin-commutes-cap-c1}
Let $(S, \delta)$ be as in Situation \ref{situation-setup}. Let $X$ be
locally of finite type over $S$. Let $(\mathcal{L}, s, i : D \to X)$
be a triple as in Definition \ref{definition-gysin-homomorphism}.
Let $\mathcal{N}$ be an invertible $\mathcal{O}_X$-module.
Then $i^*(c_1(\mathcal{N}) \cap \alpha) = c_1(i^*\mathcal{N}) \cap i^*\alpha$
in $\CH_{k - 2}(D)$ for all $\alpha \in \CH_k(X)$.
\end{lemma}

\begin{proof}
With exactly the same proof as in Lemma \ref{lemma-gysin-factors}
this follows from Lemmas
\ref{lemma-pushforward-cap-c1},
\ref{lemma-cap-commutative}, and
\ref{lemma-gysin-factors-general}.
\end{proof}

\begin{lemma}
\label{lemma-gysin-commutes-gysin}
Let $(S, \delta)$ be as in Situation \ref{situation-setup}. Let $X$ be locally
of finite type over $S$. Let $(\mathcal{L}, s, i : D \to X)$ and
$(\mathcal{L}', s', i' : D' \to X)$ be two triples as in
Definition \ref{definition-gysin-homomorphism}. Then the diagram
$$
\xymatrix{
\CH_k(X) \ar[r]_{i^*} \ar[d]_{(i')^*} & \CH_{k - 1}(D) \ar[d]^{j^*} \\
\CH_{k - 1}(D') \ar[r]^{(j')^*} & \CH_{k - 2}(D \cap D')
}
$$
commutes where each of the maps is a Gysin map.
\end{lemma}

\begin{proof}
Denote $j : D \cap D' \to D$ and $j' : D \cap D' \to D'$ the closed
immersions corresponding to $(\mathcal{L}|_{D'}, s|_{D'}$ and
$(\mathcal{L}'_D, s|_D)$. We have to show that
$(j')^*i^*\alpha = j^* (i')^*\alpha$ for all $\alpha \in \CH_k(X)$.
Let $W \subset X$ be an integral closed subscheme of dimension $k$.
Let us prove the equality in case $\alpha = [W]$. We will deduce
it from the key formula.

\medskip\noindent
We let $\sigma$ be a nonzero meromorphic section of $\mathcal{L}|_W$
which we require to be equal to $s|_W$ if $W \not \subset D$.
We let $\sigma'$ be a nonzero meromorphic section of $\mathcal{L}'|_W$
which we require to be equal to $s'|_W$ if $W \not \subset D'$.
Write
$$
\text{div}_{\mathcal{L}|_W}(\sigma) =
\sum \text{ord}_{Z_i, \mathcal{L}|_W}(\sigma)[Z_i] = \sum n_i[Z_i]
$$
and similarly
$$
\text{div}_{\mathcal{L}'|_W}(\sigma') =
\sum \text{ord}_{Z_i, \mathcal{L}'|_W}(\sigma')[Z_i] = \sum n'_i[Z_i]
$$
as in the discussion in Section \ref{section-key}.
Then we see that $Z_i \subset D$ if $n_i \not = 0$ and
$Z'_i \subset D'$ if $n'_i \not = 0$. For each $i$, let $\xi_i \in Z_i$
be the generic point. As in Section \ref{section-key} we choose
for each $i$ an element
$\sigma_i \in \mathcal{L}_{\xi_i}$, resp.\ $\sigma'_i \in \mathcal{L}'_{\xi_i}$
which generates over $B_i = \mathcal{O}_{W, \xi_i}$
and which is equal to the image of
$s$, resp.\ $s'$ if $Z_i \not \subset D$, resp.\ $Z_i \not \subset D'$.
Write $\sigma = f_i \sigma_i$ and $\sigma' = f'_i\sigma'_i$ so that
$n_i = \text{ord}_{B_i}(f_i)$ and
$n'_i = \text{ord}_{B_i}(f'_i)$.
From our definitions it follows that
$$
(j')^*i^*[W] =
\sum \text{ord}_{B_i}(f_i) \text{div}_{\mathcal{L}'|_{Z_i}}(\sigma'_i|_{Z_i})
$$
as cycles and
$$
j^*(i')^*[W] =
\sum \text{ord}_{B_i}(f'_i) \text{div}_{\mathcal{L}|_{Z_i}}(\sigma_i|_{Z_i})
$$
The key formula (Lemma \ref{lemma-key-formula}) now gives the equality
$$
\sum \left(
\text{ord}_{B_i}(f_i) \text{div}_{\mathcal{L}'|_{Z_i}}(\sigma'_i|_{Z_i}) -
\text{ord}_{B_i}(f'_i) \text{div}_{\mathcal{L}|_{Z_i}}(\sigma_i|_{Z_i})
\right) =
\sum \text{div}_{Z_i}(\partial_{B_i}(f_i, f'_i))
$$
of cycles. Note that $\text{div}_{Z_i}(\partial_{B_i}(f_i, f'_i)) = 0$ if
$Z_i \not \subset D \cap D'$ because in this case either $f_i = 1$
or $f'_i = 1$. Thus we get a rational equivalence between our specific
cycles representing $(j')^*i^*[W]$ and $j^*(i')^*[W]$ on $D \cap D' \cap W$.
By Remark \ref{remark-infinite-sums-rational-equivalences}
the result follows for general $\alpha$.
\end{proof}
















